% Chapter 6 — includable by main.tex
\chapter{Leme o pumpanju}
\label{chap:6}

\begin{quote}
\itshape
Ta stvar zove se \textbf{svojstvo pumpanja}.

Intuicija kaže "treba pamtiti neograničeno puno", ali to nije formalni dokaz.
Leme o pumpanju daju nam precizno oružje: ako jezik nema određeno
"periodično" svojstvo, onda ne može biti u toj klasi.
\end{quote}

\section{Lema o pumpanju za regularne jezike}
% ================================================================

\begin{propozicija}[Lema o pumpanju, Sipser 1.70]
Neka je $L$ regularan jezik. Tada postoji broj $p \geq 1$
(``duljina pumpanja'') takav da se svaka riječ $w \in L$ duljine
$|w| \geq p$ može zapisati kao $w = xyz$ uz uvjete:
\begin{enumerate}[label=(\roman*)]
  \item $y \neq \eps$ \quad (pumpa je neprazna),
  \item $|xy| \leq p$ \quad (pumpa je blizu početka),
  \item $xy^i z \in L$ za sve $i \geq 0$ \quad (pumpanje cuva jezik).
\end{enumerate}
\end{propozicija}

\textit{Dokaz.}
Neka KA $M$ prepoznaje $L$, i neka je $p = |Q|$ broj stanja.
Za riječ $w \in L$ s $|w| \geq p$: čitajući prvih $p$ znakova,
automat prolazi kroz $p+1$ stanja $r_0, r_1, \ldots, r_p$.
Po \textbf{Dirichletovom principu} (golubinjak), neka dva stanja
su ista: $r_i = r_j$ za neke $0 \leq i < j \leq p$.

Definirajmo:
\[
  x = w_{0:i}, \quad y = w_{i:j}, \quad z = w_{j:|w|}.
\]
Tada:
\begin{itemize}
  \item $y \neq \eps$ jer $j > i$, dakle $|y| = j - i > 0$.
  \item $|xy| = j \leq p$.
  \item Automat koji procita $xy^i z$: ulazi u $r_i$ nakon $x$,
        petlja $i$ puta kroz $r_i \to \cdots \to r_j = r_i$,
        zatim nastavlja od $r_j = r_i$ kroz $z$ i zavrsava u istom
        zavrsnom stanju kao i za $w$. \qed
\end{itemize}

\begin{kljucno}
Lemu koristimo \textbf{kontrapositivno}: ako za svaki $p$ možemo
naći riječ $w \in L$ koja se ne može pumpati (za svaki rastav
$xyz$ koji zadovoljava (i) i (ii), postoji $i$ takav da
$xy^iz \notin L$), tada $L$ nije regularan.

Dokaz je \textbf{igra}: naš protivnik odabire $p$, mi biramo $w$,
protivnik bira rastav $xyz$, mi biramo $i$. Pobjedujemo ako dobijemo
$xy^iz \notin L$.
\end{kljucno}

% ================================================================
\section{Primjena: $L_=$ nije regularan}
% ================================================================

\begin{propozicija}
Jezik $L_= = \{a^n b^n \mid n \geq 1\}$ nije regularan.
\end{propozicija}

\textit{Dokaz.}
Pretpostavimo suprotno: neka je $L_=$ regularan s duljinom pumpanja $p$.

Odaberemo riječ $w = a^p b^p \in L_=$ (duljine $2p \geq p$, dakle
zadovoljava uvjet). Prema lemi, postoji rastav $w = xyz$ s
$y \neq \eps$ i $|xy| \leq p$.

Iz $|xy| \leq p = |a^p|$ slijedi da $x$ i $y$ leze unutar bloka $a$-ova,
dakle $y = a^m$ za neki $m \geq 1$.

Pumpajmo s $i = 0$: $xy^0z = xz = a^{p-m}b^p$. Da bi ovo bilo u $L_=$,
trebamo jednak broj $a$-ova i $b$-ova: $p - m = p$, ali $m \geq 1$
pa $p - m < p$. Kontradikcija: $xz \notin L_=$. \qed

\begin{primjer}[Tipična strategija dokaza]
Struktura dokaza je uvijek ista:
\begin{enumerate}
  \item \textbf{Pretpostavimo} da je $L$ regularan s duljinom pumpanja $p$.
    \item \textbf{Odaberemo} konkretnu riječ $w \in L$ duljine $\geq p$
      (tipično $w$ ovisi o $p$, npr.\ $w = a^p b^p$).
    \item \textbf{Analiziramo} sve moguće rastave $xyz$ koji zadovoljavaju
      uvjete (i) i (ii). Često uvjet $|xy| \leq p$ snažno ograničava
      gdje može ležati $y$.
  \item \textbf{Pumpamo} s nekim $i$ (tipicno $i = 0$ ili $i = 2$) i
        pokazujemo da $xy^i z \notin L$.
  \item \textbf{Kontradikcija}: $L$ nije regularan.
\end{enumerate}
\end{primjer}

\begin{primjer}[$\{w \in \{a,b\}^* : |w|_a = |w|_b\}$ nije regularan]
Neka je $L = \{w : |w|_a = |w|_b\}$. Odaberemo $w = a^p b^p$.
Sve kao gore: $y = a^m$, pumpanje s $i = 0$ daje $a^{p-m}b^p$
s razlicitim brojem $a$-ova i $b$-ova. \qed
\end{primjer}

\begin{primjer}[$\{0^n : n \text{ je prost}\}$ nije regularan]
Neka je $L = \{0^n : n \text{ je prost}\}$.
Odaberemo $w = 0^p$ gdje je $p$ prvi prost broj $\geq p$ duljine
pumpanja (takav postoji jer ima beskonačno prostih).

Rastav: $w = xyz$, $|y| = m \geq 1$, $|xy| \leq p$.
Pumpamo s $i = |xz| + 1$: duljina od $xy^i z$ je
$|x| + i \cdot m + |z| = |xz| + im = |xz|(1 + m) + m \cdot |xz|$\ldots

Koristimo drugi trik: $|xy^i z| = |xz| + im$.
Uzmemo $i = |xz|$, tada $|xy^iz| = |xz| + |xz| \cdot m = |xz|(1+m)$
--- djeljivo s $|xz|$ i s $(1+m)$, dakle nije prost za $|xz| > 1$.
\qed
\end{primjer}

\begin{zadatak}
Dokazite da sljedeći jezici nisu regularni (koristeći lemu o pumpanju):
\begin{enumerate}[label=(\alph*)]
  \item $\{a^n b^{2n} \mid n \geq 0\}$
  \item $\{w \in \{a,b\}^* : w = w^R\}$ (palindromi)
  \item $\{1^{n^2} \mid n \geq 0\}$ (kvadrati u unarnom zapisu)
  \item $\{a^i b^j c^k \mid i + j = k\}$
\end{enumerate}
\end{zadatak}

% ================================================================
\section{Lema o pumpanju za bezkontekstne jezike}
% ================================================================

Ista ideja, ali za BK-jezike i s pet dijelova umjesto tri.

\begin{propozicija}[Lema o pumpanju za BK, Sipser 2.34]
Neka je $L$ bezkontekstan jezik. Tada postoji broj $p \geq 1$
takav da se svaka riječ $s \in L$ s $|s| \geq p$ može zapisati
kao $s = uvxyz$ uz uvjete:
\begin{enumerate}[label=(\roman*)]
  \item $vy \neq \eps$ \quad (barem jedna od pumpi je neprazna),
  \item $|vxy| \leq p$ \quad (pumpe su blizu jedna drugoj),
  \item $uv^i x y^i z \in L$ za sve $i \geq 0$ \quad (pumpanje cuva jezik).
\end{enumerate}
\end{propozicija}

\textit{Ideja dokaza.}
Neka BKG $G$ (u ChNF) generira $L$, i neka je $b = \max |desna strana|`
(za ChNF, $b = 2$). Za dovoljno dugu riječ $s$, stablo parsiranja
mora imati duboku granu --- toliko duboku da se neka varijabla $R$
mora ponoviti duž iste grane (po Dirichletovom principu, ako
grana ima više od $|V|$ čvorova).

Istaknemo zadnja dva pojavljivanja $R$ u toj grani:
\[
  S \Rightarrow^* uRz \Rightarrow^* uvRyz \Rightarrow^* uvxyz = s.
\]
Budući da $R \Rightarrow^* vRy$, možemo iterirati: $R \Rightarrow^* v^i R y^i$,
a zatim zatvoriti s $R \Rightarrow^* x$. Time dobijemo $uv^i xy^i z \in L$
za svaki $i \geq 0$.

Uvjet $|vxy| \leq p$ dolazi iz toga da biramo \emph{zadnja} dva
pojavljivanja $R$ u dubini: podstablo ispod nižeg $R$ nije predugo,
pa riječ koju ono generira nije preduga. \qed

% ================================================================
\section{Primjena: $L_\equiv$ nije bezkontekstan}
% ================================================================

\begin{propozicija}
Jezik $L_\equiv = \{a^n b^n c^n \mid n \geq 1\}$ nije bezkontekstan.
\end{propozicija}

\textit{Dokaz.}
Pretpostavimo suprotno: neka je $p$ duljina pumpanja za $L_\equiv$.

Odaberemo $s = a^p b^p c^p$ (duljine $3p \geq p$). Prema lemi,
postoji rastav $s = uvxyz$ s $vy \neq \eps$ i $|vxy| \leq p$.

Ključna opservacija: $|vxy| \leq p$, a $a$-ovi i $c$-ovi su udaljeni
$2p > p$ znakova. Dakle $vxy$ ne može sadržavati i $a$-ove i $c$-ove
istovremeno.

\textbf{Slučaj 1}: $vxy$ ne sadrzi $c$-ove.
Tada je svih $p$ slova $c$ u dijelu $z$. Pumpamo s $i = 2$:
riječ $uv^2xy^2z$ ima više $a$-ova ili $b$-ova (jer smo pumpali $v$
ili $y$ koji sadrže $a$-ove ili $b$-ove), ali i dalje $p$ slova $c$.
Dakle $uv^2xy^2z \notin L_\equiv$.

\textbf{Slučaj 2}: $vxy$ ne sadrzi $a$-ove (analogno, ili BNSO).
Svih $p$ slova $a$ je u $u$. Pumpamo s $i = 0$:
$uxz$ ima $p$ slova $a$, ali manje od $p$ slova $b$ ili $c$
(jer $vy \neq \eps$ pa smo nešto uklonili). Dakle $uxz \notin L_\equiv$. \qed

\begin{primjer}[Nezatvorenost BK na presjek]
Definiramo:
\begin{align*}
  L_{a=b} &= \{a^n b^n c^k : n, k \geq 1\} = L_= \cdot c^+, \\
  L_{b=c} &= \{a^k b^n c^n : n, k \geq 1\} = a^+ \cdot L_=.
\end{align*}
Oba su bezkontekstna (konkatenacija BK i regularnog).
Ali $L_{a=b} \cap L_{b=c} = L_\equiv$ --- a upravo smo dokazali da
$L_\equiv$ nije bezkontekstan!

Zaključak: \textbf{BK nije zatvorena na presjek}. Iz de Morgana slijedi
ni na komplement.
\end{primjer}

\begin{zadatak}
Dokazite da sljedeći jezici nisu bezkontekstni:
\begin{enumerate}[label=(\alph*)]
  \item $\{a^n b^n c^n d^n \mid n \geq 1\}$
  \item $\{ww \mid w \in \{a,b\}^*\}$ (kvadrati)
  \item $\{a^i b^j c^k \mid i < j < k\}$
  \item $\{1^{n^2} \mid n \geq 0\}$ (je li ovo BK?)
\end{enumerate}
\end{zadatak}

% ================================================================
\section{Usporedba dviju lema}
% ================================================================

\begin{center}
\renewcommand{\arraystretch}{1.6}
\begin{tabular}{lll}
\hline
& \textbf{Reg (lema za Reg)} & \textbf{BK (lema za BK)} \\
\hline
Rastav & $xyz$ (3 dijela) & $uvxyz$ (5 dijelova) \\
Pumpe & $y$ (jedna pumpa) & $v$ i $y$ (dvije pumpe) \\
Uvjet nepraznosti & $y \neq \eps$ & $vy \neq \eps$ \\
Uvjet lokalizacije & $|xy| \leq p$ & $|vxy| \leq p$ \\
Pumpanje & $xy^iz \in L$ & $uv^ixy^iz \in L$ \\
Ključni alat dokaza & Dirichlet na stanjima KA & Dirichlet na varijablama u stablu \\
\hline
\end{tabular}
\end{center}

Obje leme imaju isti oblik dokaza neregularnosti/nebezkontekstnosti:
pretpostavi, odaberi rijec, analiziraj rastav, pumpa, kontradikcija.

\begin{kljucno}
Leme o pumpanju su \textbf{nuzni} uvjeti za regularnost/bezkontekstnost,
ali ne i \textbf{dovoljni}! Postoje jezici koji zadovoljavaju lemu o
pumpanju za Reg, a nisu regularni. Lema se koristi samo za dokazivanje
\emph{neregularnosti} --- ne može dokazati regularnost.
\end{kljucno}

% ================================================================
\section*{Sazetak}
% ================================================================

\medskip\noindent
\begin{tabular}{ll}
  Lema za Reg & $w = xyz$, $y \neq \eps$, $|xy| \leq p$, $xy^iz \in L$ \\
  Lema za BK  & $s = uvxyz$, $vy \neq \eps$, $|vxy| \leq p$, $uv^ixy^iz \in L$ \\
  Strategija & pretpostavi $\to$ odaberi riječ $\to$ analiziraj $\to$ pumpa $\to$ kontradikcija \\
\end{tabular}

\medskip\noindent
U sljedećem poglavlju prelazimo na kontekstne jezike i ograničene
automate --- prvu klasu za kojom trebamo ``pravu'' traku.

